\section{Compito tre: ricerca del cammino minimo}

Il terzo compito descrive la progettazione e l'implementazione dell'algoritmo ricorsivo \texttt{CAMMINOMIN} per la ricerca del cammino minimo globale tra una cella di origine $O$ e una cella di destinazione $D$. L'algoritmo opera su un grafo implicito i cui nodi sono costituiti dalle celle della griglia e i cui archi rappresentano i cammini liberi.

\subsection{Definizione ricorsiva e flusso esecutivo}

L'algoritmo si basa sulla scomposizione del percorso ottimale in una successione di sotto-cammini liberi. A ogni livello di ricorsione:
\begin{enumerate}
    \item \textbf{Verifica del caso base}: si calcola la chiusura dell'origine corrente. Se la destinazione $D$ appartiene a tale chiusura (ossia è raggiungibile con un singolo cammino libero di tipo 1 o tipo 2), la ricorsione si interrompe immediatamente ritornando la distanza ideale $d_{\text{lib}}(O,D)$ e i landmark associati;
    \item \textbf{Espansione della frontiera}: se $D$ non è direttamente raggiungibile, si estrae la frontiera locale della chiusura. Se la frontiera è vuota, il ramo corrente rappresenta un vicolo cieco. Viene quindi ritornato un costo infinito e una sequenza vuota di landmark;
    \item \textbf{Esplorazione e ritracciamento}: per ciascun candidato di frontiera $F$ si calcola il costo parziale $l_F = d_{\text{lib}}(O,F)$. Se la mossa supera le condizioni di potatura, si effettua una chiamata ricorsiva con origine in $F$ e destinazione $D$, aggiornando dinamicamente lo stato della griglia.
\end{enumerate}

Le chiusure calcolate nei livelli superiori dell'albero di ricorsione vengono marcate come ostacoli temporanei, per prevenire cicli e per evitare di riesplorare aree già visitate.

\subsection{Analisi della potatura: riga 16 contro riga 17}

Per ridurre lo spazio di ricerca e per evitare l'esplorazione di rami non ottimali, la specifica introduce due condizioni di potatura alternative:
\begin{itemize}
    \item \textbf{Potatura debole (riga 16)}: la cella $F$ viene esplorata solo se la distanza ideale dal nodo corrente è inferiore al cammino ottimale globale provvisorio trovato finora:
    \begin{equation}
    l_F < \text{lunghezzaMin}
    \label{eq:pruning_weak}
    \end{equation}
    \item \textbf{Potatura forte (riga 17)}: sfrutta la stima ammissibile della distanza libera residua da $F$ alla destinazione $D$, potando il ramo se la somma del cammino parziale e della stima minima rimanente non può battere l'ottimo corrente:
    \begin{equation}
    l_F + d_{\text{lib}}(F, D) < \text{lunghezzaMin}
    \label{eq:pruning_strong}
    \end{equation}
\end{itemize}

La potatura forte è più onerosa per via del calcolo aggiuntivo di $d_{\text{lib}}$, ma riduce i nodi visitati in presenza di ostacoli complessi.

\subsection{Potatura branch-and-bound globale}

Lo pseudocodice valuta la condizione di riga 16/17 rispetto a \texttt{lunghezzaMin}, cioè il minimo \emph{locale} della singola invocazione. L'implementazione la rafforza confrontandola con un \textbf{limite superiore globale} $L_{\text{glob}}$, condiviso da tutte le invocazioni ricorsive e aggiornato ogni volta che si trova un cammino completo più breve. Un candidato di frontiera $F$ viene quindi esplorato solo se
\begin{equation}
\ell_{\text{acc}} + l_F + h(F) < L_{\text{glob}}
\label{eq:branch_and_bound}
\end{equation}
dove $\ell_{\text{acc}}$ è la lunghezza già accumulata dall'origine principale e $h(F)$ vale $0$ nella potatura debole (riga 16) e $d_{\text{lib}}(F,D)$ in quella forte (riga 17). Poiché $d_{\text{lib}}$ \textbf{non sovrastima mai la distanza reale}, la stima è ammissibile e la potatura non scarta mai un cammino ottimo: si tratta di una potatura branch-and-bound. Finché non si è trovato il primo cammino completo, $L_{\text{glob}}$ vale $+\infty$ e non si pota nulla. Nella configurazione predefinita la potatura non richiede alcuna passata preliminare: il limite emerge dalla ricerca stessa. È però possibile anticiparne l'attivazione con un innesco opzionale, descritto di seguito.

\subsection{Innesco goloso del limite superiore globale}

Finché $L_{\text{glob}}$ vale $+\infty$, la potatura branch-and-bound non scarta alcun ramo: l'effetto benefico descritto sopra si manifesta solo dopo che la ricerca ha trovato il primo cammino completo. Per anticiparlo, l'implementazione offre un innesco opzionale (parametro \texttt{greedy\_seed}, flag \texttt{-{}-greedy-seed}): alla sola chiamata radice viene eseguita una discesa golosa senza backtracking, che a ogni passo sceglie unicamente la cella di frontiera con $d_{\text{lib}}$ minima verso $D$ (lo stesso criterio euristico già usato per l'ordinamento della frontiera, diapositiva 66 del testo dell'elaborato) e prosegue da lì, marcando via via le chiusure attraversate come ostacoli temporanei. La discesa termina in un cammino completo qualsiasi oppure in un vicolo cieco (frontiera vuota, nel qual caso non fornisce alcun innesco); la griglia viene comunque ripristinata allo stato originale prima che la ricerca vera e propria abbia inizio.

\textbf{Correttezza}: un cammino completo qualsiasi, per definizione, non è mai più corto del cammino minimo, quindi la sua lunghezza è sempre un limite superiore valido. Usarla per inizializzare $L_{\text{glob}}$ prima dell'esplorazione non può quindi mai far scartare un cammino ottimo: è la stessa argomentazione di ammissibilità già impiegata per la potatura branch-and-bound stessa. Una sottigliezza da gestire: la condizione di potatura usa la disuguaglianza stretta, per cui un cammino di lunghezza esattamente pari a $L_{\text{glob}}$ verrebbe scartato durante l'esplorazione; se l'innesco è già ottimo, la ricerca completa non ne ritroverebbe la sequenza. Per questo l'innesco non si limita a impostare $L_{\text{glob}}$, ma inizializza anche il minimo e la sequenza restituiti dall'invocazione radice: se la ricerca completa non trova nulla di strettamente migliore, viene restituito il cammino goloso stesso.

\textbf{Legittimità rispetto alla consegna}: questo innesco è una deviazione dallo pseudocodice sullo stesso piano di quella già argomentata per il limite superiore globale, e trova fondamento negli stessi punti della traccia: la specifica ammette esplicitamente che l'implementazione non debba essere fedele allo pseudocodice purché produca gli stessi effetti operativi (diapositiva 34 del testo dell'elaborato), e discute essa stessa, a proposito delle righe 16 e 17, la possibilità di sostituire la condizione di potatura con una più efficace pur di non comprometterne la correttezza (diapositive 36--37 del testo dell'elaborato). Per questo motivo l'innesco resta disattivo in modo predefinito, cosicché il comportamento di base rimanga la traduzione letterale dello pseudocodice; l'attivazione è un'estensione facoltativa e dichiarata.

Un esperimento dedicato (\texttt{scripts/esperimento\_ottimizzazioni.py}, output in \texttt{results/greedy\_seed\_results.json} e pannello sinistro della figura~\ref{fig:ottimizzazioni}) confronta chiamate ricorsive e tempo con e senza innesco, per entrambe le potature, su tre istanze di difficoltà crescente; l'esito è istruttivo proprio perché il beneficio misurato è nullo o modesto, per ragioni diverse a seconda dell'istanza. Si noti che, per le istanze che raggiungono il tempo limite, il numero di chiamate ricorsive registrato è quello scandito dall'orologio di sistema entro il tempo concesso, non un valore deterministico: può quindi variare leggermente da un'esecuzione all'altra a parità di macchina, oltre che fra macchine diverse:
\begin{itemize}
    \item \textbf{Sulla griglia $18\times18$} (la stessa della figura~\ref{fig:esplorazione_potatura}) la discesa golosa riesce e trova subito il cammino di lunghezza $24.63$, che è già l'ottimo: ma la potatura forte, priva di innesco, risolve la stessa istanza in sole $2$ chiamate ricorsive contro l'unica chiamata necessaria con l'innesco, quindi il margine di miglioramento possibile era comunque minimo;
    \item \textbf{Sulla griglia $40\times40$} dell'esempio di cammino complesso (figura~\ref{fig:cammino_risolto_complesso}) la discesa golosa \textbf{fallisce}: si esaurisce in un vicolo cieco senza mai raggiungere $D$, e non fornisce alcun innesco. La ricerca con e senza \texttt{-{}-greedy-seed} produce quindi risultati identici ($11\,303$ chiamate con potatura forte, $60.43$ di lunghezza), perché l'innesco è di fatto assente. Questo è un risultato negativo interessante di per sé: mostra concretamente perché una discesa golosa senza ritracciamento non è sufficiente in generale a risolvere il problema, e perché il ritracciamento completo di \texttt{CAMMINOMIN} resta necessario;
    \item \textbf{Sulla griglia $100\times100$} (tempo limite di $60$ s per ciascuna corsa) né la potatura debole né quella forte completano la ricerca, con o senza innesco: la lunghezza del miglior cammino trovato entro il tempo limite coincide comunque con o senza innesco ($174.02$ con potatura debole, $150.45$ con potatura forte), segno che l'innesco non ha fornito in questa esecuzione un limite iniziale abbastanza stringente da spostare l'esito qualitativo. Il numero di chiamate ricorsive entro il tempo limite varia invece misurabilmente: con potatura debole scende da $1\,504\,225$ a $1\,341\,857$ (circa l'$11\%$ in meno), con potatura forte da $300\,998$ a $299\,428$ (circa lo $0{,}5\%$ in meno) — una riduzione reale ma comunque insufficiente a cambiare l'esito entro il tempo concesso.
\end{itemize}
Nel complesso, sulle istanze provate l'innesco goloso non produce benefici pratici decisivi: sulle istanze facili la potatura forte con ordinamento euristico converge già in pochissime chiamate; su quelle difficili la discesa senza ritracciamento tende a incontrare un vicolo cieco proprio dove servirebbe di più, oppure il vantaggio in chiamate ricorsive resta troppo piccolo per cambiare l'esito qualitativo entro il tempo limite. L'estensione resta comunque corretta, non peggiora mai il risultato (§4.4, argomento di ammissibilità) e disponibile come opzione (\texttt{-{}-greedy-seed}), documentata per completezza della valutazione critica degli esiti sperimentali richiesta dalla specifica (diapositiva 70 del testo dell'elaborato).

Per produrre questo confronto, \texttt{camminomin()} accetta un parametro opzionale \texttt{visited\_closures}: se fornito, ogni chiamata ricorsiva vi accumula le celle della propria chiusura, tracciando così l'impronta complessiva della ricerca a fini di analisi. Il parametro non altera in alcun modo la logica dell'algoritmo, non è utilizzato quando omesso (valore predefinito \texttt{None}) ed è impiegato solo dallo script che genera la figura seguente.

Per rendere tangibile l'effetto della potatura, la figura~\ref{fig:esplorazione_potatura} confronta, sulla stessa coppia origine--destinazione di una griglia $18 \times 18$ con ostacoli ad agglomerati, l'unione di tutte le chiusure esplorate durante l'intera ricerca nei due casi. La \textbf{potatura debole} (riga 16) allaga quasi l'intero spazio libero raggiungibile ($260$ celle su $324$), perché continua a inseguire rami la cui unica colpa è avere percorso finora più strada, anche quando sono geometricamente lontani dalla destinazione. La \textbf{potatura forte} (riga 17) traccia invece un corridoio molto più mirato verso $D$ ($97$ celle), poiché la stima $d_{\text{lib}}(F,D)$ scarta subito i rami che, nel migliore dei casi, non potrebbero comunque competere con il minimo corrente. Su questa istanza il tempo di calcolo passa da $21$ ms a meno di un millisecondo.

\begin{figure}[H]
\centering
\includegraphics[width=\textwidth]{esplorazione_potatura_debole_vs_forte.png}
\caption{Impronta di esplorazione (unione delle chiusure visitate) a confronto fra potatura debole e forte, stessa coppia origine--destinazione.}
\label{fig:esplorazione_potatura}
\end{figure}

\subsection{Controllo preliminare di raggiungibilità con i componenti connessi}

La ricerca ricorsiva ha un punto debole strutturale sulle coppie irraggiungibili: quando la destinazione si trova in una regione isolata da ostacoli (per esempio all'interno di un recinto chiuso), l'algoritmo non ha modo di accorgersene se non esaurendo l'intera frontiera di ogni ramo, con un costo che cresce in modo esponenziale prima di poter rispondere che il cammino non esiste.

Per eliminare questo caso patologico l'implementazione esegue, soltanto alla chiamata radice, un riempimento a inondazione iterativo (visita in ampiezza con una coda, funzione \texttt{compute\_components()}) che etichetta ogni cella libera con l'identificativo del proprio componente connesso rispetto al vicinato a otto celle, in tempo e spazio $O(R \cdot C)$. Se origine e destinazione ricadono in componenti diversi, la funzione ritorna subito distanza infinita senza avviare alcuna ricorsione. La correttezza discende dalla specifica stessa (diapositiva 3 del testo dell'elaborato): il passo diagonale fra due celle libere è sempre ammesso, anche quando due ostacoli si toccano per lo spigolo, quindi la connettività a otto coincide esattamente con la raggiungibilità e il controllo non può mai dichiarare irraggiungibile una coppia che non lo sia. L'esito del calcolo resta identico in entrambi i casi (distanza infinita), cambia soltanto il costo per arrivarci.

L'effetto è quantificato da un esperimento dedicato (\texttt{scripts/esperimento\_ottimizzazioni.py}, output in \texttt{results/component\_check\_results.json} e pannello destro della figura~\ref{fig:ottimizzazioni}): su una griglia $50 \times 50$ con densità $0{,}2$ e un recinto chiuso, con destinazione interna al recinto e origine esterna, il controllo preliminare (sola etichettatura dei componenti connessi, $O(R\cdot C)$) risponde in $0{,}006$ secondi con $0$ chiamate ricorsive, mentre la stessa interrogazione senza controllo raggiunge il tempo limite di $600$ secondi dopo oltre $27$ milioni di invocazioni ricorsive senza essere riuscita a dimostrare l'irraggiungibilità: quattro ordini di grandezza di differenza fra i due casi. Il controllo è attivo in modo predefinito e disattivabile dalla riga di comando con \texttt{-{}-no-component-check} per riprodurre il confronto. Va osservato che il beneficio riguarda soltanto le coppie realmente disconnesse: sulle coppie raggiungibili il controllo aggiunge un costo lineare trascurabile e la ricerca procede come prima.

\begin{figure}[H]
\centering
\includegraphics[width=0.95\textwidth]{7_ottimizzazioni.png}
\caption{Sinistra: invocazioni ricorsive con e senza innesco goloso, per potatura debole/forte, sulle tre istanze di prova. Destra: invocazioni ricorsive con e senza controllo dei componenti connessi sulla coppia irraggiungibile.}
\label{fig:ottimizzazioni}
\end{figure}

\subsection{Ottimizzazioni fisiche e implementative}

L'implementazione prevede anche le seguenti ottimizzazioni strutturali:
\begin{itemize}
    \item \textbf{Ordinamento euristico della frontiera}: prima di scorrere il ciclo sulle celle della frontiera, queste vengono ordinate per valori crescenti della distanza libera dalla destinazione tramite l'ordinamento nativo di Python (Timsort, $O(n \log n)$)~\cite{python-sort}. Questa strategia esplora per primi i percorsi più promettenti, riducendo il numero di rami da visitare prima di aggiornare il minimo globale;
    \item \textbf{Ritracciamento sul posto}: realizzato marcando le celle della chiusura corrente nella matrice condivisa con un valore costante (ostacolo temporaneo) ed eseguendo il ripristino in tempo costante durante la risalita, senza allocazione di memoria.
\end{itemize}

Un tentativo iniziale aveva introdotto anche la memoizzazione dei sotto-problemi, con chiave data da origine, destinazione e insieme delle celle temporaneamente bloccate. Il tentativo ha rivelato però un difetto logico: il risultato di ogni sotto-problema dipende, attraverso la potatura globale, dalla lunghezza accumulata delle chiamate esterne, che non fa parte della chiave. Un risultato memorizzato con una lunghezza accumulata elevata, e quindi con i rami ottimali già potati, poteva essere reimpiegato erroneamente in una chiamata con lunghezza accumulata minore, producendo una risposta subottimale e violando la simmetria richiesta dalla specifica (diapositiva 64 del testo dell'elaborato). La memoizzazione è stata pertanto rimossa, mentre le prestazioni restano garantite dalla potatura branch-and-bound globale e dall'ordinamento euristico della frontiera.

Tutte queste ottimizzazioni sono implementate come parametri booleani opzionali di \texttt{camminomin()} (\texttt{use\_strong\_pruning}, \texttt{randomize\_frontier}).

\newpage
\subsection{Requisiti funzionali implementati}

Il software prevede le seguenti funzionalità richieste dalla specifica:
\begin{itemize}
    \item \textbf{Ricostruzione del cammino fisico (\texttt{reconstruct\_path})}: a partire dalla sequenza di landmark ritornata ($\langle (O,0), (F_1, t_1), \dots, (D, t_D) \rangle$), il sistema ricostruisce l'esatta sequenza continua di celle attraversabili (mosse cardinali e diagonali) unendo i singoli tratti liberi di tipo 1 o tipo 2 e validandone l'attraversabilità geometrica;
    \item \textbf{Gestione dell'interruzione per superamento del tempo limite}: l'ispezione del tempo trascorso a ogni chiamata ricorsiva, consente di arrestare la ricerca in modo sicuro se si supera il tempo limite programmato dall'utente; salvando e restituendo la sequenza di celle di lunghezza minore individuata fino a quel momento.
\end{itemize}

La figura~\ref{fig:cammino_risolto} mostra un primo esempio didattico di cammino minimo ricostruito: in azzurro trasparente il percorso fisico completo, con la spezzata e i marcatori romboidali che ne evidenziano la sequenza di landmark (visualizzazione ibrida, diapositiva 53 del testo dell'elaborato). Su questa istanza la destinazione si trova quasi in linea diagonale con l'origine, per cui la sequenza si riduce a un solo landmark intermedio.

\begin{figure}[H]
\centering
\includegraphics[width=0.62\textwidth]{cammino_risolto.png}
\caption{Cammino minimo ricostruito su un'istanza semplice: pochi ostacoli, sequenza breve di landmark.}
\label{fig:cammino_risolto}
\end{figure}

La figura~\ref{fig:cammino_risolto_complesso} riporta invece un caso più rappresentativo del comportamento dell'algoritmo su griglie realistiche: una griglia $40 \times 40$ con agglomerati, barre e recinti (densità $0.25$), dove il cammino minimo fra i due angoli opposti richiede $7$ landmark intermedi per aggirare gli ostacoli via via incontrati. Ogni deviazione visibile nella spezzata corrisponde a una cella di frontiera da cui l'algoritmo ha aperto un nuovo cammino libero verso la destinazione.

\begin{figure}[H]
\centering
\includegraphics[width=0.72\textwidth]{cammino_risolto_complesso.png}
\caption{Cammino minimo su una griglia $40 \times 40$ densamente ostacolata: $7$ landmark intermedi.}
\label{fig:cammino_risolto_complesso}
\end{figure}
